%%%%%%%%%%%%%%%%%%%%%%%%%%%%%%%%%%%%%%%%%%%%%%%%%%%%%%%%%%%%%%%%%%%%%%%%
% Portada oscura: identificación institucional, título, superficie
% |1/(1+z^2)| coloreada por la fase (los dos polos z = ±i que limitan el
% radio de convergencia de la serie real) y datos de autoría.
%%%%%%%%%%%%%%%%%%%%%%%%%%%%%%%%%%%%%%%%%%%%%%%%%%%%%%%%%%%%%%%%%%%%%%%%
\begin{titlepage}
\thispagestyle{empty}
\begin{tikzpicture}[remember picture, overlay]
  \coordinate (O) at (current page.south west);
  \coordinate (C) at ($(O)+(10.5cm,0)$);

  % Logotipo y cabecera institucional
  \node[anchor=north, inner sep=0] at ($(C)+(0,27.6cm)$)
    {\includegraphics[height=1.45cm]{logo_usal_blanco.png}};
  \draw[grismedio, line width=0.35pt] ($(O)+(3.0cm,25.6cm)$) -- ($(O)+(18.0cm,25.6cm)$);
  \node[anchor=north, align=center, inner sep=0] at ($(C)+(0,25.15cm)$)
    {{\color{gris}\scshape\small\textls[140]{\Facultad\ \ ·\ \ \Titulacion}}};

  % Asignatura, tema y título
  \node[anchor=north, align=center, inner sep=0, text width=16cm] at ($(C)+(0,23.3cm)$)
    {{\color{cian}\scshape\large\textls[180]{\Asignatura}}\\[6pt]
     {\color{gris}\scshape\normalsize\textls[140]{\NumTema}}\\[20pt]
     {\color{blanco}\fontsize{27}{34}\selectfont\bfseries\TituloTema}\\[14pt]
     {\color{cian}\rule{3.2cm}{0.7pt}}\\[12pt]
     {\color{texto}\itshape\large\SubtituloTema}};

  % Superficie |f(z)| de f(z) = 1/(1+z^2), coloreada por arg f(z)
  \node[anchor=center, inner sep=0] at ($(C)+(0,10.9cm)$) {%
    \begin{tikzpicture}
      % La fase salta de +180 a -180 grados sobre el eje imaginario |y|>1:
      % se dibuja la superficie en dos mitades (primero la lejana, x>0) para que el
      % sombreado interpolado no recorra todo el círculo cromático en el corte.
      \begin{axis}[superficie, width=15.5cm, height=11.5cm,
          view={-28}{24}, axis lines=none, clip=false,
          y domain=-2:2, samples=\resPx, samples y=\resPy,
          zmin=0, zmax=4.6, point meta min=-180, point meta max=180]
        \addplot3[surf, z buffer=sort, domain=0.0001:2,
            point meta={-atan2(2*x*y,1+x^2-y^2)}]
          {min(1/sqrt((1+x^2-y^2)^2+4*x^2*y^2+0.00005),4.6)};
        \addplot3[surf, z buffer=sort, domain=-2:-0.0001,
            point meta={-atan2(2*x*y,1+x^2-y^2)}]
          {min(1/sqrt((1+x^2-y^2)^2+4*x^2*y^2+0.00005),4.6)};
      \end{axis}
    \end{tikzpicture}};

  % Epígrafe: la serie entera y la cota del radio
  \node[anchor=north, align=center, inner sep=0] at ($(C)+(0,5.5cm)$)
    {\color{texto}$\displaystyle
      y(x)=\sum_{n=0}^{\infty}a_n\,(x-x_0)^n,\qquad
      \frac{1}{1+x^{2}}=\sum_{n=0}^{\infty}(-1)^n x^{2n}\quad(\abs{x}<1)$};

  % Autoría, lugar y curso
  \draw[grismedio, line width=0.35pt] ($(O)+(3.0cm,4.2cm)$) -- ($(O)+(18.0cm,4.2cm)$);
  \node[anchor=north west, inner sep=0, align=left] at ($(O)+(3.0cm,3.75cm)$)
    {{\scshape\small\color{cian}\textls[100]{Autor}}\\[3pt]
     {\large\color{blanco}\miNombre}};
  \node[anchor=north, inner sep=0, align=center] at ($(C)+(0,3.75cm)$)
    {{\scshape\small\color{cian}\textls[100]{Profesorado}}\\[3pt]
     {\small\color{texto}\miProfesorado}};
  \node[anchor=north east, inner sep=0, align=right] at ($(O)+(18.0cm,3.75cm)$)
    {{\scshape\small\color{cian}\textls[100]{\tipotrabajo}}\\[3pt]
     {\color{gris}\miUbicacion}};
\end{tikzpicture}
\end{titlepage}
